\section{\color{cmured} \textbf{Things in TEST --- results}}

\begin{itemize}
    \item \textbf{ACT-c10 on dp01} --- trained (15k steps, loss 0.213) and mid-screen.
    So far 0/4 strict seats but \textbf{1 physically completed insertion}
    (plug left in jack, 109 frames parked $>$14\,mm aligned): the gentle 0.1\,mm
    floor alone already restores enough fine alignment to finish, even with zero
    recovery data.

    \item \textbf{40-episode eval matrix} --- \textbf{done} (08:42--12:34).
    First-ever closed-loop seats: ACT-depaused 1/20 (on a \emph{held-out}
    condition), pi05-depaused 1/20 (train condition), plus 4 stable deep
    insertions 1--2\,mm short. Revised the ``always drifts'' N=1 verdict:
    depaused models are low-rate succeeders ($\sim$5\% strict), and the
    bottleneck moved from ``can't find the hole'' to the final millimetres.

    \item \textbf{ACT-c10 + tier-2 recovery} --- \textbf{the decisive experiment,
    and the thesis held.} Two variants trained: recovery-\emph{only}
    (weak alone: 0/6, one bore entry --- it lacks full approaches, as expected)
    and \textbf{combined} (depaused base + recovery, 1{,}562 eps):
    \textbf{2/6 physically completed insertions incl.\ one held-out} ---
    plug fully in the jack, hand withdrawn --- scored 0/6 only because the
    seat criterion requires the plug \emph{grasped} at depth (the gripper slips
    during the terminal push). Recovery data is confirmed as the missing
    ingredient; the new (much smaller) bottleneck is terminal grip slip.
    Follow-up staged: consistent 0.1\,mm-floor mix
    (\texttt{dp01} base + \texttt{tier2rec\_dp01}, pushing to HF now) to test
    whether preserved slow-push behavior converts ``insert-then-slip'' into held
    seats.

    \item {\color{cmured}low priority} \textbf{pi05 on dp01 + tier2} --- not
    started (team/cluster). Now the recommended production retrain given the
    ACT results; all datasets staged on HF.
\end{itemize}

\subsection*{\color{cmured} New datasets behind these tests}

\paragraph{De-paused demonstrations.}
The scripted expert nearly halts at its pre-dock standoff ($\sim$0.02\,mm/frame,
19\,mm before the socket), so a third of all frames carried near-zero action labels
--- which $\pi_{0.5}$ learned as an absorbing ``stop here'' state. We resample each
episode by arc length (keep a frame only after $\geq$0.2\,mm of motion) and recompute
all action labels from the retained poses. This removes 23\% of frames and cuts the
slow-frame fraction from 27.1\% to 0.1\%, leaving recovery segments and the action
distribution untouched.

\paragraph{Finer variant (dp01).}
The 0.2\,mm floor overshot: de-paused models move decisively but lose the expert's
careful final alignment. A 0.1\,mm floor (10.8\% cut) removes parking while
preserving the fine-alignment regime.

\paragraph{Recovery demonstrations (tier-2).}
No demonstration showed how to return to the path after an error --- the failure every
model shared. Two remedies (562 eps / 50{,}755 frames, same plant and DR as the base
corpus): \emph{off-nominal starts} (200 eps beginning 15--22\,mm out and up to
16\,mm off-axis, so every frame is corrective), and \emph{perturbation recovery}
(a brief 2--6\,mm lateral hijack of the command, then expert correction). Since labels
come from pose differences, the shove would read as intentional motion --- episodes
are therefore split at the perturbation, keeping only the recovery. A 0.1\,mm re-trim
accompanies the dp01 base.

\footnotesize{HuggingFace: \texttt{cable\_drmix\_*\_\{depaused,dp01\}},
\texttt{cable\_tier2rec\_*\_\{depaused,dp01\}} (Parallax-Worlds).}

\subsection*{\color{cmured} If time permits, things to try}

\begin{enumerate}
    \item \textbf{RL-based teacher.} Replace/augment the scripted expert with an RL
    policy trained in sim (privileged state, force-aware reward): it would discover
    terminal-push and slip-avoidance behavior the hand-tuned controller lacks, and
    its natural state-visitation already covers off-path recovery --- no kick
    machinery needed.
    \item \textbf{DAgger-style targeted recovery.} Roll out the \emph{current} policy,
    detect divergence, and relabel its own visited states with the scripted expert.
    All infrastructure exists (eval harness + controller); this replaces synthetic
    kicks with data at exactly the states the policy actually reaches.
    \item \textbf{Temporal ensembling for ACT} (\texttt{temporal\_ensemble\_coeff}),
    currently disabled. Averages overlapping chunk predictions at eval time --- zero
    retraining, historically reduces jerk and chunk-boundary drift.
    \item \textbf{Wrist-image dropout during training.} Attention maps show the wrist
    view carries the task while the front view anchors on background; randomly
    blanking the wrist image would force front-view competence --- the sensor needed
    for recovery once the plug leaves the wrist frame.
    \item \textbf{Observation history} (\texttt{n\_obs\_steps} $>$ 1). A single-frame
    policy cannot distinguish ``drifting left'' from ``statically offset left'';
    velocity context should sharpen corrections.
    \item \textbf{Terminal-phase grip fix.} Diagnose the insert-then-slip failure:
    higher grasp force / friction in the plant, slower commanded push in the last
    2\,mm (dp01 may already fix this), or score released-but-seated as success.
    \item \textbf{Recovery-share sweep.} Combined used $\sim$28\% recovery data;
    try 40--50\% before generating anything new.
    \item \textbf{Real-bench fine-tune.} Record supervised bench trials at 60\,Hz;
    even tens of real episodes (especially corrections) may close most of the
    sim-to-real gap.
\end{enumerate}

\paragraph{Where the policy looks.}
To probe \emph{how} the policies use their two views, we visualize the ACT decoder's
cross-attention --- head-averaged weights of the ten action queries over the encoder's
image tokens, mapped back onto each camera --- across timesteps of both nominal and
recovery episodes (Fig.~\ref{fig:attn}). The division of labor is consistent: the
wrist camera carries the task (52--56\% of attention mass, concentrated in a sharp
hotspot on the plug tip and the plug--socket contact zone), while the front camera
attends mostly to static background structure, apparently serving as a coarse pose
anchor rather than a servoing signal. This asymmetry offers a mechanistic account of
the universal no-recovery failure: once a policy drifts far enough that the plug and
socket leave the wrist view, its primary task sensor is gone, and the front view was
never trained to take over. Notably, on off-path frames from the recovery dataset the
front camera's attention visibly shifts onto the jack --- suggesting recovery training
recruits exactly the missing front-view competence, and giving us a cheap,
rollout-free diagnostic for future checkpoints: front-camera attention on the socket
in off-path states should co-vary with recovery ability. (Attention is suggestive
rather than causal; the model has a single decoder layer, so these maps summarize its
only cross-attention stage.)
